\subsection{Caso 9. Priorización y correlación de alertas con IA y SIEM}
\label{caso:siem-alertas}

\begin{situacion}
El centro de monitoreo recibe múltiples alertas que deben agruparse y priorizarse.
\end{situacion}

\paragraph{Procedimiento.}
\begin{itemize}[leftmargin=*, topsep=2pt, itemsep=2pt]
	\item Normalizar eventos de firewall, servidor, Zeek/Suricata y sistema operativo; eliminar duplicados sin perder trazabilidad.
	\item Definir características agregadas por host y ventana temporal, y probar un detector de anomalías con datos controlados.
	\item Usar un LLM únicamente para resumir alertas anonimizadas y sugerir el triage; verificar toda recomendación contra la evidencia.
\end{itemize}

\paragraph{Indicadores que deben buscarse.}
Alertas correlacionadas por host/tiempo, puntuación de riesgo explicable y señales convergentes en varias fuentes.

\paragraph{Mitigación y defensa.}
Revisión humana antes de acciones disruptivas; control de acceso a los registros; evaluación de sesgo y deriva; registro de decisiones y posibilidad de reversión.

\paragraph{Evidencias que se entregan.}
Diagrama del pipeline, ejemplo de alerta enriquecida, criterios de prioridad y registro de validación humana.

\paragraph{Preguntas de análisis.}
\begin{enumerate}[leftmargin=*, topsep=2pt, itemsep=2pt]
	\item ¿Cómo explicar una alerta a un analista?
	\item ¿Qué riesgo supone automatizar bloqueos basados solo en una puntuación?
\end{enumerate}

\begin{table}[H]
	\caption{Fuentes de eventos del pipeline de correlación}
	\label{tab:fuentes-siem}
	\centering
	\footnotesize
	\renewcommand{\arraystretch}{1.25}
	\begin{tabularx}{\textwidth}{@{}>{\raggedright\arraybackslash}p{0.22\textwidth} Y >{\raggedright\arraybackslash}p{0.16\textwidth} >{\raggedright\arraybackslash}p{0.18\textwidth}@{}}
		\toprule
		\textbf{Fuente} & \textbf{Tipo de evento} & \textbf{Formato} & \textbf{Responsable} \\
		\midrule
		Firewall & & & \\
		Servidor web & & & \\
		Zeek / Suricata & & & \\
		Sistema operativo & & & \\
		SIEM / Wazuh & & & \\
		& & & \\
		\bottomrule
	\end{tabularx}
\end{table}

\registrohallazgo{Caso 9}
